\documentclass[11pt,a4paper]{article}
\usepackage[utf8]{inputenc}
\usepackage[T1]{fontenc}
\usepackage{lmodern}
\usepackage[french]{babel}
\usepackage{amsmath,amssymb,mathtools}
\usepackage{icomma}
\usepackage{xcolor}
\usepackage[most]{tcolorbox}
\usepackage{enumitem}
\usepackage[a4paper,margin=2.2cm,headheight=26pt,headsep=10pt]{geometry}
\usepackage{fancyhdr}
\usepackage[hidelinks]{hyperref}
\definecolor{primary}{RGB}{222,49,99}
\definecolor{primarydark}{RGB}{150,22,55}
\tcbset{boxbase/.style={enhanced,breakable,boxrule=0.9pt,arc=2.5pt,
  left=3mm,right=3mm,top=2.3mm,bottom=2.3mm,fonttitle=\bfseries,coltitle=white}}
\newtcolorbox[auto counter]{exo}[1][]{boxbase,colback=primary!4,colframe=primary,
  title={\textsc{Exercice}~\thetcbcounter\ifx\relax#1\relax\relax\else\textmd{ --- \itshape #1}\fi}}
\newcommand{\R}{\mathbb{R}}
\newcommand{\ds}{\displaystyle}
\pagestyle{fancy}\fancyhf{}
\fancyhead[L]{\small\textcolor{primarydark}{\textbf{Thème 2} — Logarithmes}}
\fancyhead[R]{\small\textcolor{primarydark}{\textsc{Facile}}}
\fancyfoot[C]{\small\thepage}
\renewcommand{\headrulewidth}{0.8pt}
\begin{document}

\begin{center}
{\LARGE\bfseries\color{primarydark} Niveau Facile — Exercices}\\[2pt]
{\color{primary}\rule{0.45\textwidth}{1.2pt}}\\[3pt]
{\small\itshape Une seule mécanique à la fois. Aucune calculatrice.}
\end{center}
\vspace{1mm}

\begin{exo}[la définition : un log est un exposant]
Calcule (utilise $\log_a x=y\iff x=a^{y}$) :
\begin{enumerate}[label=\alph*),leftmargin=2em,topsep=1pt,itemsep=1pt]
  \item $\log_2 8$ \qquad
  \item $\log_3 81$ \qquad
  \item $\log_{10} 1000$ \qquad
  \item $\log_5 1$ \qquad
  \item $\ln e$
\end{enumerate}
\end{exo}

\begin{exo}[développer avec les trois propriétés]
Développe (produit $\to$ somme, quotient $\to$ différence, puissance $\to$ produit), $x,y>0$ :
\begin{enumerate}[label=\alph*),leftmargin=2em,topsep=1pt,itemsep=1pt]
  \item $\log_a(xy)$ \qquad
  \item $\ds\log_a\Bigl(\frac{x}{y}\Bigr)$ \qquad
  \item $\log_a\bigl(x^{3}\bigr)$ \qquad
  \item $\ln\bigl(x^{2}\bigr)$ \qquad
  \item $\log_a\bigl(x^{2}y\bigr)$
\end{enumerate}
\end{exo}

\begin{exo}[valeurs remarquables]
Donne la valeur exacte :
\begin{enumerate}[label=\alph*),leftmargin=2em,topsep=1pt,itemsep=1pt]
  \item $\log_{10} 10$ \qquad
  \item $\ln e^{4}$ \qquad
  \item $\log_2 1$ \qquad
  \item $\ds\ln\frac{1}{e}$ \qquad
  \item $\log_{10} 0{,}01$
\end{enumerate}
\end{exo}

\begin{exo}[appliquer $\ln$ pour isoler un exposant]
Résous dans $\R$ en appliquant $\ln$ (ou $\log$) aux deux membres :
\begin{enumerate}[label=\alph*),leftmargin=2em,topsep=1pt,itemsep=1pt]
  \item $e^{x}=7$ \qquad
  \item $10^{x}=100$ \qquad
  \item $e^{2x}=e^{5}$ \qquad
  \item $e^{-x}=3$ \qquad
  \item $10^{x}=5$
\end{enumerate}
\end{exo}

\begin{exo}[le piège $\ln^{2}x$ contre $\ln x^{2}$]
Pour $x>0$, réponds par \emph{vrai} ou \emph{faux} et corrige si nécessaire :
\begin{enumerate}[label=\alph*),leftmargin=2em,topsep=1pt,itemsep=1pt]
  \item $\ln(x^{2})=2\ln x$. \qquad
  \item $\ln^{2}x=2\ln x$. \qquad
  \item $(\ln x)^{2}=\ln^{2}x$. \qquad
  \item $\ln(x^{2})=\ln^{2}x$.
\end{enumerate}
\end{exo}

\end{document}
